\section*{Introduction}

Small uncrewed aircraft are moving from demonstration flights to routine parcel logistics. In the United States, the Federal Aviation Administration has proposed a performance-based framework that would normalise beyond-visual-line-of-sight package delivery and replace case-by-case waivers\cite{faa2025part108}, while European airspace is being organised around U-space services for dense low-altitude traffic\cite{sesar2023uspace,kopardekar2016utm}. Life-cycle and routing studies indicate that drone delivery can reduce energy use and emissions relative to ground vans for light parcels\cite{stolaroff2018energy,kirschstein2020comparison,chiang2019impact}, and a large operations-research literature addresses the assignment and routing layer of such services\cite{dorling2017vrp,macrina2020drone,cheng2020energy,bi2024truck}. The operational bottleneck, however, is increasingly the flight layer: a fleet of delivery aircraft must fly through a cluttered urban district in turbulent wind, with degraded satellite positioning, among other members of its own fleet and non-cooperative aircraft, and with a battery whose usable energy depends on payload, ambient conditions and the health of the propulsion system. The safety case of such an operation is dominated by rare events, a single collision or a single forced landing after energy depletion, rather than by the mean delivery time.

Reinforcement learning (RL) has become a credible way of synthesising flight policies in this regime. Deep RL controllers have reached champion level in drone racing\cite{kaufmann2023champion,song2023reaching}, learned adaptive flight in strong winds\cite{oconnell2022neuralfly,zhang2025extreme} and produced time-optimal end-to-end quadrotor control\cite{ferede2024endtoend}, while multi-agent RL has been used for tactical separation assurance in structured and unstructured airspace\cite{brittain2021jais,brittain2022scalable,chen2024uam,sharifi2026hetero}. Off-policy actor-critic methods such as the soft actor-critic (SAC)\cite{haarnoja2018sac} and TD3\cite{fujimoto2018td3}, and on-policy methods such as PPO\cite{schulman2017ppo,schulman2016gae} and its parameter-shared multi-agent form\cite{yu2022mappo}, are the usual workhorses. All of them maximise an expected discounted return. An expected-value criterion is indifferent between a policy that is always slightly late and a policy that is usually early but occasionally catastrophic, which is precisely the distinction a delivery operator must make. Domain randomisation\cite{tobin2017domain} widens the training distribution but does not by itself change what the policy is optimised for.

Distributional RL offers the missing representation. Instead of estimating the expected return, the critic estimates the full distribution of returns, either as a categorical distribution\cite{bellemare2017distributional} or through quantiles\cite{dabney2018qr,dabney2018iqn}, with truncated quantile critics controlling overestimation in continuous control\cite{kuznetsov2020tqc} and distributional soft actor-critic variants supporting risk-sensitive objectives\cite{ma2025dsac}; the theory is now consolidated\cite{bellemare2023book}, and a distributional code for value has even been reported in the dopamine system\cite{dabney2020dopamine}. Once a return distribution is available, a coherent tail-risk measure such as the conditional value-at-risk (CVaR)\cite{rockafellar2000cvar} can replace the expectation. CVaR-optimal control has a well-developed theory in Markov decision processes\cite{chow2015cvar,bauerle2011avar,wang2023minimax} and policy-gradient\cite{tamar2015cvar,tang2019wcpg} and constrained\cite{yang2021wcsac} algorithmic forms, and tail-risk measures are increasingly regarded as the natural language for robot safety\cite{majumdar2020risk,akella2025riskaware,chapman2022risk,hakobyan2019risk,royset2025riskadaptive}. Two difficulties are, however, well documented. First, a fixed risk level is a blunt instrument: a low level makes the agent timid everywhere and a high level makes it indifferent to hazards, and risk-averse optimisation can become blind to successful strategies unless the risk level is annealed\cite{greenberg2022cesor}. Second, greedy CVaR maximisation of the return-to-go is time-inconsistent, so the risk level should in principle evolve along the trajectory\cite{lim2022distributional,bauerle2011avar}.

These difficulties have motivated risk-sensitive navigation agents with conditioned or adapted risk levels. The risk-conditioned distributional soft actor-critic learns a single policy for a continuum of risk levels\cite{choi2021rcdsac}; adaptive risk-tendency methods set the CVaR level of a nano-drone from the spread of its learned return distribution\cite{liu2023art,seres2023flight}; automatic risk adaptation and online epistemic risk adaptation choose the level from novelty or from an online total-variation criterion\cite{schubert2021ara,wu2026drlora}; and recent marine and ground-robot studies tie the level to measured uncertainty or to obstacle distance\cite{huang2025usv,wang2026diqn,lin2023usv,zhang2025driqn,wang2025adaptive,schneider2024quadruped}. In cooperative multi-agent RL, CVaR has been used to train risk-sensitive value factorisations\cite{qiu2021rmix,shen2023riskq}. Across this body of work the adaptation signal is either internal to the learner (return spread, novelty, epistemic uncertainty) or a single geometric quantity, and the risk measure is applied to the total return, in which efficiency variability and catastrophic events are mixed. None of these methods reasons about the energy state of the vehicle, although for battery-powered delivery aircraft energy depletion is as catastrophic as a collision and is driven by the same exogenous disturbances.

Energy feasibility has so far been treated separately, through energy models and fixed reserves. Rotary-wing power models derived from blade-element momentum theory\cite{zeng2019energy} underpin most delivery-energy studies, yet comparisons across models show large discrepancies\cite{zhang2021energy}, and mission-planning studies emphasise that wind can turn a feasible plan into an infeasible one\cite{thibbotuwawa2020weather,cheng2020energy}. Deployed return-to-home logic typically keeps a fixed multiplicative margin over a nominal estimate, which is either wasteful or unsafe as soon as conditions drift. Conformal prediction supplies distribution-free, finite-sample uncertainty bounds\cite{vovk2005algorithmic,angelopoulos2023gentle} and has been adopted for safe planning, perception-based control and decision making in robotics\cite{lindemann2023safe,lindemann2025formal,yang2023perception,lekeufack2024cdt}; its adaptive and online forms retain long-run coverage under arbitrary distribution shift\cite{gibbs2021aci,angelopoulos2023pid,barber2023beyond} and have been used for motion planning among dynamic agents and crowd navigation\cite{dixit2023adaptive,sheng2024pomdp,yao2025crowd}. Safe RL more broadly relies on constraints, shields and barrier functions\cite{garcia2015survey,brunke2022safe,gu2024review,achiam2017cpo,alshiekh2018shielding,ames2019cbf}, and backup-policy safety filters are receiving renewed attention\cite{kim2026backup}. To the best of our knowledge, no prior study couples a risk-sensitive multi-agent flight policy with a statistically calibrated energy reserve, and none evaluates both under controlled, family-by-family distribution shift of wind, sensing, traffic and energy.
